
\documentclass{article}

\usepackage[margin=1in]{geometry}
\usepackage{amsmath, amssymb, amsthm, mathtools}
\usepackage{enumitem}
\usepackage{hyperref}

\newcommand{\R}{\mathbb{R}}
\newcommand{\Prob}{\mathbb{P}}
\newcommand{\E}{\mathbb{E}}
\newcommand{\norm}[1]{\| #1 \|}
\newcommand{\tr}{\operatorname{tr}}
\newcommand{\rank}{\operatorname{rank}}
\newcommand{\diag}{\operatorname{diag}}
\newcommand{\clip}{\operatorname{clip}}
\newcommand{\bA}{\mathbf{A}}
\newcommand{\bB}{\mathbf{B}}
\newcommand{\bN}{\mathbf{N}}
\newcommand{\bM}{\mathbf{M}}
\newcommand{\bX}{\mathbf{X}}
\newcommand{\bU}{\mathbf{U}}
\newcommand{\bV}{\mathbf{V}}
\newcommand{\bS}{\mathbf{S}}
\newcommand{\bQ}{\mathbf{Q}}

\theoremstyle{plain}
\newtheorem{theorem}{Theorem}
\newtheorem{lemma}{Lemma}
\newtheorem{corollary}{Corollary}
\newtheorem{proposition}{Proposition}

\theoremstyle{definition}
\newtheorem{algorithm}{Algorithm}

\title{Robust Low-Rank Approximation via Column-Norm Sampling}
\author{}
\date{\today}

\begin{document}
\maketitle

Let $\bB=\bA+\bN\in\R^{n\times n}$, where $\bA\succeq 0$ is an unknown positive semidefinite matrix and $\bN$ is an unknown perturbation. Bakshi, Chepurko, and Woodruff~\cite{BCW21} study robust low-rank approximation in this model under the assumptions
\[
    \norm{\bN}_F^2\leq \eta\norm{\bA}_F^2
    \qquad\text{and}\qquad
    \norm{\bN_{i,*}}_2^2\leq c\norm{\bA_{i,*}}_2^2
    \quad\text{for all }i\in[n],
\]
where $c$ is a fixed constant. Their algorithm is parameterized by a measure of diagonal corruption:
\[
    \phi_{\max}:=\max_{i\in[n]}\frac{\bA_{ii}}{|\bB_{ii}|}.
\]
Their Theorem 5.11 gives a rank-$k$ approximation $\bX$ satisfying
\[
    \norm{\bA-\bX}_F^2
    \leq
    \norm{\bA-\bA_k}_F^2+(\varepsilon+\sqrt{\eta})\norm{\bA}_F^2
\]
using $\widetilde O(\phi_{\max}^2nk/\varepsilon)$ queries to $\bB$.

\medskip
% In this note, we show how to sample rows and columns from 

% Throughout, assume
% \[
%     \norm{\bN}_F^2\leq \eta\norm{\bA}_F^2
% \]
% and suppose that we are given a number $\Phi\geq 1$ such that
% \[
%     \bA_{ii}\leq \Phi | \bB_{ii}|
%     \qquad\text{for every }i\in[n].
% \]
% In the notation of Bakshi--Chepurko--Woodruff, one may take $\Phi=\phi_{\max}$ and $\rho=\sqrt{\eta}$.

% \section{Clipping}
Let $\Phi$ be a number satisfying $\Phi \geq \max\{1, \phi_{\max}\}$. We first define the clipped matrix $\widetilde{\bB}$ entrywise by
\[
    \widetilde{\bB}_{ij}
    :=
    \clip\!\left(
        \bB_{ij},
        -\Phi\sqrt{|\bB_{ii}|\,|\bB_{jj}|},
        \Phi\sqrt{|\bB_{ii}|\,|\bB_{jj}|}
    \right),
\]
where
\[
    \clip(x,-\tau,\tau):=\min\{\tau,\max\{-\tau,x\}\}.
\]
A query to $\widetilde{\bB}_{ij}$ is very efficient when the diagonal entries of $\bB$ are known. Bakshi, Chepurko, and Woodruff state that we can sample from $\widetilde \bB$ without increasing corruption to $\bA$~\cite{BCW21}. The following lemma formalizes this statement:

\begin{lemma}[Clipping lemma]
The clipped matrix satisfies
\[
    \widetilde{\bB}_{ij}^{\,2}
    \leq
    \Phi^2|\bB_{ii}|\,|\bB_{jj}|
    \qquad\text{for all }i,j\in[n].
\]
Moreover, we may write
\[\widetilde\bB = \bA + \widetilde\bN\]
for some matrix $\widetilde\bN$ satisfying $\|\widetilde\bN\|^2_F\leq \|\bN\|^2_F$. 
\end{lemma}

\begin{proof}
The first claim follows directly from the definition of clipping. For the
second claim, we first note
\[
    \bA_{ii}\leq \Phi|\bB_{ii}|
    \qquad\text{and}\qquad
    \bA_{jj}\leq \Phi|\bB_{jj}|.
\]
by the definition of $\phi_{\max}$. Moreover, $\bA_{ij}^2\leq \bA_{ii}\bA_{jj}$ since $\bA$ is psd, so
\[
    |\bA_{ij}|
    \leq
    \Phi\sqrt{|\bB_{ii}|\,|\bB_{jj}|}.
\]
So the clipping interval for entry $ij$ contains 
$\bA_{ij}$. In particular,
\[
    |\bB_{ij}-\bA_{ij}|
    =
    |\bB_{ij}-\widetilde{\bB}_{ij}|
    +
    |\widetilde{\bB}_{ij}-\bA_{ij}|
    \geq
    |\widetilde{\bB}_{ij}-\bA_{ij}|.
\]
Therefore
\[
    (\widetilde{\bB}_{ij}-\bA_{ij})^2
    \leq
    \left(\bB_{ij}-\bA_{ij}\right)^2
    =
    \bN_{ij}^2.
\]
Summing over all entries gives $\|\widetilde\bB - \bA\|^2_F\leq \|\bN\|^2_F$, which is the desired result.
\end{proof}

Given entry-wise access to $\bB$, we consider the following procedure to sample indices with probability proportional to the row (column) norms of $\widetilde\bB$. Let
\[
    D:=\sum_{i=1}^n |\bB_{ii}|.
\]
Assume $D>0$. In one trial, 
\begin{enumerate}
    \item Sample indices $i,j$ independently from $\frac{\diag(|\bB|)}{D}$.
    \item Accept $i$ with probability
    $\frac{\widetilde{\bB}_{ij}^{\,2}}{\Phi^2|\bB_{ii}|\,|\bB_{jj}|}.$
\end{enumerate}

\begin{lemma}
Conditioned on acceptance,
\[
    p_i
    =
    \frac{\|\widetilde{\bB}_{i,*}\|_2^2}{\|\widetilde{\bB}\|_F^2}.
\]
Similarly, if the acceptance probability is changed to
\[
    \frac{\widetilde{\bB}_{ji}^{\,2}}
    {\Phi^2|\bB_{ii}|\,|\bB_{jj}|},
\]
then
\[
    p_i
    =
    \frac{\|\widetilde{\bB}_{*,i}\|_2^2}{\|\widetilde{\bB}\|_F^2}.
\]
In either case, the expected number of trials until acceptance is at most
$\Phi^2 n$.
\end{lemma}

\begin{proof}
The acceptance probability is well-defined because the clipping lemma gives
\[
    \widetilde{\bB}_{ij}^{\,2}
    \leq
    \Phi^2|\bB_{ii}|\,|\bB_{jj}|.
\]
If $|\bB_{ii}|\,|\bB_{jj}|=0$, then the pair $(i,j)$ is sampled with probability
zero, so the value of the acceptance rule on that pair is irrelevant.

In one trial, the probability of outputting $i$ is
\[
    \frac{|\bB_{ii}|}{D}
    \sum_{j=1}^n
    \frac{|\bB_{jj}|}{D}
    \cdot
    \frac{\widetilde{\bB}_{ij}^{\,2}}
    {\Phi^2|\bB_{ii}|\,|\bB_{jj}|}
    =
    \frac{\norm{\widetilde{\bB}_{i,*}}_2^2}
    {\phi_{\max}^2 D^2}.
\]
Summing over all indices gives total acceptance probability
\[
    \frac{\norm{\widetilde{\bB}}_F^2}{\Phi^2 D^2}
\]
and conditioning on acceptance, the probability of outputting $i$ is
\[
    \frac{\norm{\widetilde{\bB}_{i,*}}_2^2}{\norm{\widetilde{\bB}}_F^2}.
\]
The column-norm statement is identical after replacing
$\widetilde{\bB}_{ij}$ by $\widetilde{\bB}_{ji}$.

It remains to bound the acceptance probability. Since
$\Phi\geq 1$, the diagonal entries are not clipped, so
\[
    D^2
    =
    \left(\sum_{i=1}^n |\bB_{ii}|\right)^2
    \leq
    n\sum_{i=1}^n \bB_{ii}^2
    =
    n\sum_{i=1}^n \widetilde{\bB}_{ii}^{\,2}
    \leq
    n\norm{\widetilde{\bB}}_F^2.
\]
Thus the acceptance probability satisfies
\[
    \frac{\norm{\widetilde{\bB}}_F^2}{\Phi^2D^2}
    \geq
    \frac{1}{\Phi^2 n}.
\]

\end{proof}


\section{Robust Low-Rank Approximation}
We use the following result of Frieze, Kannan, and Vempala \cite{FKV04}. Let $\bM\in\R^{n\times n}$. Suppose one can sample rows of $\bM$ with
probability exactly
\[
    p_i=\frac{\norm{\bM_{i,*}}_2^2}{\norm{\bM}_F^2}
\]
Then, using $O(k/\varepsilon)$ sampled rows, there
is a randomized algorithm that outputs a rank-$k$ matrix $\bX$ satisfying
\[
    \norm{\bM-\bX}_F^2
    \leq
    \norm{\bM-\bM_k}_F^2+\varepsilon\norm{\bM}_F^2
\]
with constant probability.

\begin{theorem}
Let $\bM = \widetilde\bB$. Then the algorithm above outputs a rank-$k$ matrix $\bX$
such that, with constant probability,
\[
    \norm{\bA-\bX}_F^2
    \leq
    \norm{\bA-\bA_k}_F^2
    +
    O(\varepsilon+\sqrt{\eta})\norm{\bA}_F^2.
\]
\end{theorem}

\begin{proof}
By mimicking equations 2.12 and 2.13 of \cite{BCW21}.
% First define
% \[
%     \widetilde \bN=\widetilde{\bB}-\bA.
% \]
% By the clipping lemma,
% \[
%     \norm{\widetilde \bN}_F\leq \norm{\bN}_F\leq \sqrt{\eta}\norm{\bA}_F.
% \]
% Applying the guarantee of \cite{FKV04} to $\widetilde{\bB}$ gives
% \[
%     \norm{\widetilde{\bB}-\bX}_F^2
%     \leq
%     \norm{\widetilde{\bB}-\widetilde{\bB}_k}_F^2
%     +
%     \varepsilon\norm{\widetilde{\bB}}_F^2.
% \]
% Since $\bA_k$ is rank-$k$,
% \[
%     \norm{\widetilde{\bB}-\widetilde{\bB}_k}_F^2
%     \leq
%     \norm{\widetilde{\bB}-\bA_k}_F^2
%     =
%     \norm{\bA-\bA_k+\widetilde\bN}_F^2.
% \]
% Using Cauchy--Schwarz and $\norm{\bA-\bA_k}_F\leq\norm{\bA}_F$,
% \[
%     \norm{\bA-\bA_k+\widetilde\bN}_F^2
%     \leq
%     \norm{\bA-\bA_k}_F^2
%     +
%     O(\sqrt{\eta})\norm{\bA}_F^2.
% \]
% Also,
% \[
%     \norm{\widetilde{\bB}}_F
%     \leq
%     \norm{\bA}_F+\norm{\widetilde\bN}_F
%     \leq
%     (1+\sqrt{\eta})\norm{\bA}_F,
% \]
% so the guarantee implies
% \[
%     \norm{\widetilde{\bB}-\bX}_F^2
%     \leq
%     \norm{\bA-\bA_k}_F^2
%     +
%     O(\varepsilon+\sqrt{\eta})\norm{\bA}_F^2.
% \]
% Finally,
% \[
%     \norm{\bA-\bX}_F^2
%     =
%     \norm{\widetilde{\bB}-\bX-\widetilde\bN}_F^2.
% \]
% Since $\norm{\widetilde\bN}_F\leq\sqrt{\eta}\norm{\bA}_F$ and the preceding display also
% bounds $\norm{\widetilde{\bB}-\bX}_F=O(\norm{\bA}_F)$, another application of
% Cauchy--Schwarz gives
% \[
%     \norm{\bA-\bX}_F^2
%     \leq
%     \norm{\widetilde{\bB}-\bX}_F^2
%     +
%     O(\sqrt{\eta})\norm{\bA}_F^2.
% \]
% Combining the last two inequalities proves
% \[
%     \norm{\bA-\bX}_F^2
%     \leq
%     \norm{\bA-\bA_k}_F^2
%     +
%     O(\varepsilon+\sqrt{\eta})\norm{\bA}_F^2.
% \]
\end{proof}





\begin{thebibliography}{9}

\bibitem{BCW21}
Ainesh Bakshi, Nadiia Chepurko, and David P. Woodruff.
\newblock Robust and sample optimal algorithms for PSD low-rank approximation.
\newblock \emph{arXiv:1912.04177v5}, 2021.

\bibitem{FKV04}
Alan Frieze, Ravi Kannan, and Santosh Vempala.
\newblock Fast Monte-Carlo algorithms for finding low-rank approximations.
\newblock \emph{Journal of the ACM}, 51(6):1025--1041, 2004.

\end{thebibliography}

\end{document}
